\فصل{معماری و روش پیشنهادی}

\قسمت{مقدمه}
هدف از این فصل، تشریح دقیق معماری فنی و جزئیات پیاده‌سازی سامانه‌ی پاسخ‌گویی هوشمند طراحی شده است. این سامانه با بهره‌گیری از معماری تولید افزوده با بازیابی\LTRfootnote{Retrieval-Augmented Generation (RAG)}~\مرجع{gao2024retrieval}، تلاش می‌کند تا محدودیت‌های مدل‌های زبانی بزرگ\LTRfootnote{Large Language Models (LLMs)} در زمینه‌ی دانش ایستا و پدیده‌ی توهم\LTRfootnote{Hallucination} را از طریق اتصال به داده‌های بلادرنگ وب برطرف نماید. 

بر اساس دسته‌بندی‌های نوین در ادبیات این حوزه، معماری این سامانه در دسته‌ی \lr{RAG} پیشرفته و مدولار\LTRfootnote{Advanced / Modular RAG}~\مرجع{gao2024retrieval} قرار می‌گیرد؛ زیرا علاوه بر یک فرآیند بازیابی ساده، از خط‌لوله‌های جانبی نظیر مسیریابی هوشمند، بسط پرس‌وجو و ارجاع‌دهی دقیق بهره می‌برد. در ادامه، اجزای مختلف سیستم، از مدیریت بافتار تا خط لوله‌ی پیشرفته‌ی بازیابی و بازرتبه‌بندی، به صورت مفصل شرح داده شده است.


\قسمت{معماری کلی سامانه}
معماری پیشنهادی بر پایه‌ی یک ساختار «عامل‌گونه»\LTRfootnote{Agentic} طراحی شده است که قادر است بر اساس نوع پرسش کاربر، بین دو حالت عملیاتی مختلف سوئیچ نماید. این رویکرد را می‌توان یک معماری مبتنی بر مسیریابی\LTRfootnote{Router-based Architecture} دانست که هدف آن ایجاد تعادل میان کیفیت پاسخ‌دهی، هزینه‌ی محاسباتی و میزان تأخیر\LTRfootnote{Latency} است.

\زیرقسمت{منطق تصمیم‌گیر برای سوئیچ حالت‌ها}
سیستم در هر درخواست کاربر، بر اساس یک پارامتر کنترلی (\lr{\texttt{is\_web\_search}}) تصمیم می‌گیرد که از کدام مسیر پاسخ‌دهی استفاده کند:

\شروع{فقرات}

\فقره
\مهم{حالت پاسخ‌گویی عادی\LTRfootnote{Standard Mode}}: در این حالت، سیستم صرفاً به دانش پارامتریک\LTRfootnote{Parametric Knowledge} داخلی مدل زبانی اتکا می‌کند. این مسیر برای پرسش‌های خلاقانه، تحلیل‌های کلی یا گفت‌وگوهای ساده که نیازمند داده‌های لحظه‌ای نیستند، بهینه شده است. در این مسیر، تنها پرسش کاربر و تاریخچه‌ی گفت‌وگو به مدل ارسال می‌شود که منجر به پاسخی با کمترین تأخیر و هزینه می‌گردد.

\فقره
\مهم{حالت جست‌وجوی وب\LTRfootnote{Web Search Mode}}: زمانی که کاربر نیاز به اطلاعات به‌روز، فکت‌های دقیق یا منابع خارجی دارد، این حالت فعال می‌شود. در این مسیر، سیستم به جای پاسخ مستقیم، وارد یک خط لوله‌ی\LTRfootnote{Pipeline} پیچیده از «جست‌وجو $\rightarrow$ بازیابی $\rightarrow$ بازرتبه‌بندی $\rightarrow$ تولید» می‌گردد~\مرجع{gao2024retrieval}.

\پایان{فقرات}

این تفکیک باعث می‌شود که سیستم منابع محاسباتی سنگین‌ترِ جست‌وجوی وب و فراخوانی رابط‌های برنامه‌نویسی (\lr{API}) خارجی را تنها در موارد ضروری به کار بگیرد.


\قسمت{مدیریت تاریخچه و بافتار}
یکی از چالش‌های اصلی در سامانه‌های گفت‌وگو-محور، حفظ پیوستگی در تعاملات چندنوبتی\LTRfootnote{Multi-turn Conversations} است. برای حل این مسئله، یک سیستم مدیریت بافتار\LTRfootnote{Context Management} در لایه‌ی پایگاه‌داده\LTRfootnote{Database} و سرویس‌های منطقی پیاده‌سازی شده است.

\زیرقسمت{ذخیره‌سازی و بازیابی تاریخچه}
تمام پیام‌های تبادل شده بین کاربر و دستیار در جداول پایگاه‌داده (\lr{\texttt{Messages}} و \lr{\texttt{Chats}}) ذخیره می‌شوند. هر پیام دارای یک برچسب نقش (\lr{\texttt{role}}) است که مشخص می‌کند پیام توسط کاربر (\lr{\texttt{user}}) یا دستیار (\lr{\texttt{assistant}}) تولید شده است. هنگام ارسال هر درخواست جدید، سیستم تمام پیام‌های مربوط به آن شناسه‌ی یکتا (\lr{\texttt{chat\_id}}) را به ترتیب زمانی بازیابی می‌کند. 

\زیرقسمت{تزریق بافتار به مدل}
تاریخچه‌ی بازیابی‌شده به صورت یک رشته‌ی متنی ساختاریافته به مدل زبانی تزریق می‌شود. این کار باعث می‌شود مدل بتواند:

\شروع{فقرات}
\فقره \مهم{رفع ابهام\LTRfootnote{Disambiguation}}: ضمایری مانند «آن» یا «این» را بر اساس پیام‌های قبلی شناسایی کند.
\فقره \مهم{پیوستگی منطقی}: پاسخ‌های جدید را بر اساس جریان گفت‌وگو تنظیم کند و از تکرار مطالب قبلی بپرهیزد.
\پایان{فقرات}

شایان ذکر است که اگرچه در نسخه‌ی فعلی سامانه تمام تاریخچه‌ی گفت‌وگو بازیابی و ارسال می‌شود، در توسعه‌های آتی می‌توان از تکنیک‌های فشرده‌سازی انتخابی بافتار\LTRfootnote{Selective Context Compression}~\مرجع{wang2025optimizing} بهره برد تا ضمن کاهش مصرف توکن‌ها، تنها نوبت‌های مرتبط با پرسش نهایی به مدل تزریق شوند.


\قسمت{خط لوله‌ی جست‌وجوی وب}
این بخش (خط لوله‌ی جست‌وجوی وب\LTRfootnote{Web Search Pipeline})، قلب تپنده‌ی سامانه است و از یک ساختار چندمرحله‌ای برای تبدیل یک پرسش مبهم کاربر به یک پاسخ مستند و دقیق استفاده می‌کند.

\زیرقسمت{ماژول بسط پرس‌وجو}
به دلیل آنکه پرسش‌های کاربران لزوماً برای موتورهای جست‌وجو بهینه نیستند، سیستم از ماژول بسط پرس‌وجو\LTRfootnote{Query Expansion} برای افزایش میزان فراخوانی\LTRfootnote{Recall} استفاده می‌کند. بر اساس چارچوب‌های نوینی نظیر \lr{RQ-RAG}~\مرجع{chan2024rq}، نادیده گرفتن ابهام در پرسش‌های پیچیده می‌تواند به افت کیفیت بازیابی منجر شود.

در این مرحله، پرسش کاربر به همراه تاریخچه‌ی گفت‌وگو به یک مدل زبانی ارسال می‌شود تا با استفاده از یک دستور (پرامپت) تخصصی، پرسش اصلی را از طریق بازنویسی و تجزیه، به ۳ تا ۵ کوئری\LTRfootnote{Query} بهینه و متنوع تبدیل کند. خروجی این مرحله در قالب یک ساختار \lr{JSON} است که شامل نسخه‌های مختلفی از پرسش (از نظر واژگان و زاویه‌ی دید) می‌باشد. این کار باعث می‌شود که اگر یک کوئری نتایج ضعیفی داشت، کوئری‌های دیگر بتوانند منابع طلایی را پیدا کنند.

\زیرقسمت{لایه‌ی بازیابی اولیه}
در لایه‌ی بازیابی اولیه، کوئری‌های تولیدشده به طور موازی به رابط برنامه‌نویسی موتور جست‌وجوی تاویلی\LTRfootnote{Tavily} ارسال می‌شوند. این موتور جست‌وجو که برای هوش مصنوعی بهینه‌سازی شده است، علاوه بر لینک‌ها، محتوای خام و پالایش‌شده (\lr{\texttt{raw\_content}}) صفحات وب را برمی‌گرداند. 

این لایه از تکنیک‌های تشخیص و حذف کدهای زائد قالب\LTRfootnote{Boilerplate Removal}~\مرجع{kohlschutter2010boilerplate} استفاده می‌کند تا نویزهایی مانند منوهای سایت، تبلیغات و پانویس‌های سایت\LTRfootnote{Footers} را حذف کرده و تنها متن اصلی مقاله یا صفحه را استخراج کند. این امر منجر به کاهش چشمگیر مصرف توکن‌ها در مرحله‌ی نهایی و جلوگیری از ورود داده‌های غیرمرتبط به مدل می‌شود.

\زیرقسمت{لایه‌ی بازرتبه‌بندی ترکیبی}
در لایه‌ی بازرتبه‌بندی ترکیبی\LTRfootnote{Hybrid Reranking}، نتایج بازیابی‌شده از وب که معمولاً بر اساس شباهت‌های سطحی (مدل‌های رمزگذار دوگانه\LTRfootnote{Bi-Encoder}) رتبه‌بندی شده‌اند، مورد ارزیابی مجدد قرار می‌گیرند. برای ارتقای معیار دقت\LTRfootnote{Precision}، این لایه با استفاده از مدل‌های رمزگذار متقاطع\LTRfootnote{Cross-Encoder} پیاده‌سازی شده است~\مرجع{reimers2019sentence}.

در این سامانه، از دو مدل پیشرو یعنی کوهیر\LTRfootnote{Cohere Rerank} و جینا\LTRfootnote{Jina Reranker} استفاده شده است. برای بهینه‌سازی نهایی و افزایش پایداری سیستم، نتایج این دو مدل با استفاده از متد تلفیق رتبه متقابل\LTRfootnote{Reciprocal Rank Fusion (RRF)}~\مرجع{cormack2009reciprocal} ترکیب می‌شوند:

\شروع{فقرات}
\فقره هر دو مدل (کوهیر و جینا) به طور مستقل به اسناد رتبه می‌دهند.
\فقره سیستم به جای میانگین‌گیری ساده از امتیازات، امتیاز نهایی هر سند را از مجموع معکوسِ رتبه‌های کسب‌شده‌ی آن در هر دو مدل محاسبه می‌کند.
\فقره اسنادی که در هر دو مدل رتبه‌ی بالایی کسب کرده‌اند، به شدت تشویق شده و به عنوان مرتبط‌ترین منابع شناسایی می‌شوند و سپس به لایه‌ی تولید پاسخ ارسال می‌گردند.
\پایان{فقرات}

این رویکرد باعث می‌شود خطاهای احتمالی هر یک از مدل‌ها توسط مدل دیگر جبران شده و سیستم در مواجهه با انواع مختلف محتوا پایدارتر عمل کند.

\زیرقسمت{تولید پاسخ نهایی و سیستم ارجاع‌دهی}
در مرحله‌ی نهایی، متون پالایش‌شده و رتبه‌بندی‌شده به همراه پرسش کاربر و تاریخچه، در قالب یک پرامپت جامع به مدل زبانی ارسال می‌شود. 

سیستم از یک دستورالعمل سخت‌گیرانه (\lr{\texttt{WEB\_SEARCH\_SYSTEM\_PROMPT}}) استفاده می‌کند که مدل را مجبور می‌کند:

\شروع{فقرات}
\فقره پاسخ را \مهم{صرفاً} بر اساس متون ارائه‌شده تولید کند (تولید مستند\LTRfootnote{Grounded Generation}).
\فقره هر ادعای فکت‌محور را با یک ارجاع شماره‌دار (مانند \lr{[1], [2]}) دقیقاً به منبع مربوطه متصل کند.
\پایان{فقرات}

این رویکرد که در متون علمی به عنوان «تولید متن ارجاع‌دار»\LTRfootnote{Attributed Text Generation (ATG)} شناخته می‌شود~\مرجع{xia2024ground}، بررسی‌پذیری پاسخ‌ها را به شدت افزایش می‌دهد. در نهایت، این ارجاعات در پایگاه‌داده ذخیره شده و در رابط کاربری به صورت لینک‌های قابل کلیک نمایش داده می‌شوند تا کاربر بتواند مستقیماً به منبع اصلی مراجعه کند.


\قسمت{جزئیات فنی پیاده‌سازی}
برای دستیابی به کارایی بالا و قابلیت گسترش، از پشته‌ی فناوری\LTRfootnote{Tech Stack} زیر استفاده شده است:

\شروع{فقرات}
\فقره \مهم{زبان برنامه‌نویسی}: پایتون ۳.۱۰+ \LTRfootnote{Python 3.10+} (به دلیل غنای کتابخانه‌ها در حوزه‌ی هوش مصنوعی و پردازش داده).
\فقره \مهم{فریم‌ورک بک‌اند\LTRfootnote{Backend Framework}}: فست‌اِی‌پی‌آی\LTRfootnote{FastAPI} (به دلیل پشتیبانی از \lr{\texttt{async/await}} که برای فراخوانی‌های موازی در سرویس‌های بازیابی و بازرتبه‌بندی حیاتی است).
\فقره \مهم{مدل زبانی \lr{(LLM)}}: \lr{Gemini 2.5 Flash} (به دلیل پنجره‌ی بافتار بسیار بزرگ، سرعت بالا و توانمندی عالی در استدلال‌های مبتنی بر ارجاع)~\مرجع{google2024gemini}.
\فقره \مهم{پایگاه‌داده}: اس‌کیوال‌لایت\LTRfootnote{SQLite} با استفاده از ابزار \lr{SQLAlchemy} (برای مدیریت بهینه‌ی داده‌های کاربران، چت‌ها و ارجاعات).
\فقره \مهم{رابط کاربری (فرانت‌اند\LTRfootnote{Frontend})}: تکنولوژی‌های \lr{HTML5}، جاوا اسکریپت خالص\LTRfootnote{Vanilla JavaScript} و \lr{Tailwind CSS v4} (برای ایجاد یک محیط مدرن، سریع و بدون نیاز به فریم‌ورک‌های سنگین).
\پایان{فقرات}


\قسمت{مزایا و محدودیت‌های روش پیشنهادی}

\مهم{مزایا:}

\شروع{فقرات}
\فقره \مهم{دقت بالا و کاهش توهم}: به دلیل استفاده از لایه‌ی بازرتبه‌بندی (\LTRfootnote{Cross-Encoder}) و اجبار مدل به ارجاع‌دهی، احتمال تولید اطلاعات ساختگی به شدت کاهش یافته است.
\فقره \مهم{به‌روز بودن مطلق}: سیستم به جای اتکا به دانش قدیمی مدل، در لحظه از وب اطلاعات استخراج می‌کند.
\فقره \مهم{شفافیت و قابلیت بازبینی}: کاربر می‌تواند هر بخش از پاسخ را با کلیک بر روی ارجاعات، در منبع اصلی بررسی کند.
\پایان{فقرات}

\vspace{0.5cm} 

\مهم{محدودیت‌ها:}

\شروع{فقرات}
\فقره \مهم{تأخیر در پاسخ‌دهی (\LTRfootnote{Latency})}: به دلیل ماهیت چندمرحله‌ای سیستم (تولید کوئری $\rightarrow$ بازیابی از وب $\rightarrow$ بازرتبه‌بندی $\rightarrow$ تولید نهایی)، زمان پاسخ‌دهی در حالت جست‌وجوی وب بیشتر از حالت عادی است و ممکن است برای کاربردهای فوق‌سریع مناسب نباشد.
\فقره \مهم{وابستگی به سرویس‌های خارجی}: عملکرد بهینه‌ی سیستم مستقیماً به پایداری و در دسترس بودن وب‌سرویس‌های تاویلی، کوهیر و گوگل وابسته است.
\فقره \مهم{هزینه‌ی عملیاتی}: هر درخواست در حالت وب، اعتبار چندین وب‌سرویس را به طور همزمان مصرف می‌کند که در استقرار در مقیاس‌های بزرگ تجاری نیاز به پیاده‌سازی لایه‌های کنترل هزینه و کشینگ\LTRfootnote{Caching} خواهد داشت.
\پایان{فقرات}